\cotefile{78}
\scanpage{47}
trouver une notion convenable \inserted{de ``\unsure{répère|repère}''} ``d'\uline{épinglage}''
pour les $\mathrm{SR}n\mathrm{S}$, permettant de \deleted{les} ramener
à la question d'une structure épinglée, qu'on
pourrait soumettre à un calcul\ldots{} La théorie
de descente devrait alors permettre, \deleted{en}
\deleted{\unsure{depuis|effet|définissant}} comme le ``module'' a disparu,
de \deleted{descendre} trouver un objet épinglé
canonique sur le corps premier ($\mathbb{Q}$ ou
un $\mathbb{F}_p$), dont les autres se déduisent
par ``torsion'' par un \deleted{$(\unsure{G_k|G-k})$ to} $G_k$-torseur.
C'est sans doute ainsi dans le cas \inserted{exceptionnel} $G=\mathrm{Aff}(S)\simeq
(\mathbb{Z}/n\mathbb{Z})^*.(\mathbb{Z}/n\mathbb{Z})$, mais douteux dans le
cas $G=\mathrm{Aff}^{\pm}(S)$. Essentiellement, la question
revient à celle-ci : les $\mathrm{SR}n\mathrm{S}$ \deleted{qui} forment
une gerbe sur le corps premier $\mathbb{F}_p$
(si \deleted{\unsure{en}} $p\nmid n$, ou si $n=p\ne\unsure{2|p}$, \deleted{\unsure{mais si|premier}} ou si
$n=2p$), \deleted{\unsure{ou bien est-ce}} \inserted{celle-ci a-t-elle} une \uline{section} i.e.
existe-t-il un $\mathrm{SR}n\mathrm{S}$ sur le corps
premier $\mathbb{F}_p$ ? (Ici, bien sûr, on s'intéresse
aux structures non néc\textsuperscript{t}. scindées, où
\deleted{\unsure{les}} $S$ est un sous-schéma du fibré affine
associé à $E$, lequel sous-schéma peut
être sans points). Le mieux qu'on
% archive sheet number: 46
